\section{An elementary relative carrier with the log square factor}
\label{sec:elementary}
At integer $q=N-m$ and $M=M(m)$,
\begin{equation}\label{eq:log-square}
 \log\frac{T(N,m)}{M^q/q!}
 =\sum_{j=0}^{q-1}\log(1+j/M)
 =\frac{q(q-1)}{2M}+O(q^3/M^2).
\end{equation}
When $m\sim2N/\log N$, the main correction is asymptotic to $(\log N)^2/4$, whereas $q^3/M^2=O((\log N)^4/N)=o(1)$. The naive Poissonized cell count therefore misses a factor that grows without bound. An $O(\log^2N)$ error in an exponent cannot prove a relative equivalent.

There is nevertheless a simple implicit parametrization that retains the necessary factor. For $N>0$, let $u=u_N$ be the unique positive solution of
\begin{equation}\label{eq:ueq}
 u(u+2)e^u=2N,
\end{equation}
and put
\begin{equation}\label{eq:P}
 P(u)=u^2+4u+2.
\end{equation}
The left side of \eqref{eq:ueq} is strictly increasing from zero to infinity; also $u\sim\log N$.

\begin{theorem}[Elementary count carrier]\label{thm:elementary}
As $N\to\infty$,
\begin{equation}\label{eq:bcarrier}
 b_N=\sqrt{\frac2{P(u)}}
 \exp\!\left(\frac{Nu(u+1)}{u+2}+\frac u2+\frac{u^2}4\right)
 \left(1+O\!\left(\frac{(\log N)^4}{N}\right)\right).
\end{equation}
The implicit $u$ is retained exactly in this formula.
\end{theorem}
\begin{proof}
Set $x_0=2N/(u+2)=ue^u$ and $q_0=N-x_0=ux_0/2$. Equation \eqref{eq:ueq} gives
$\log(x_0^2/(2q_0))=u$, so $H_N'(x_0)=0$. Direct differentiation shows
\begin{equation}\label{eq:element-curvature}
 A_0=-H_N''(x_0)=\frac{u+4+2/u}{x_0},\qquad A_0q_0=P(u)/2.
\end{equation}
The differentiated estimates in Section~\ref{sec:localization} give
\[
 f_N'(x_0)=O(u^3/N),\qquad
 f_N''(x_0)=-A_0(1+O(u^2/N)).
\]
The saddle displacement is therefore $r-x_0=O(u)$, and Taylor's theorem with the uniform curvature bounds gives
\begin{equation}\label{eq:transport-small}
 f_N(r)-f_N(x_0)=O(u^4/N),\qquad A/A_0=1+O(u^2/N).
\end{equation}
For the value, use the uniform gamma quotient and Stirling expansions
\begin{align*}
 \log\Gamma(M+q)-\log\Gamma(M)
    &=q\log M+\frac{q(q-1)}{2M}+O(q^3/M^2),\\
 \log\Gamma(q+1)&=(q+1/2)\log q-q+\tfrac12\log(2\pi)+O(q^{-1}).
\end{align*}
At $x=x_0$, write $M=x_0^2(1+1/x_0)/2$. Then
$q_0\log(1+1/x_0)=u/2+O(u^2/N)$ and
$q_0(q_0-1)/(2M)=u^2/4+O(u^3/N)$, so
\[
 f_N(x_0)=q_0(u+1)+u/2+u^2/4
                -\tfrac12\log(2\pi q_0)+O(u^4/N).
\]
Apply Theorem~\ref{thm:main} with $K=1$, use \eqref{eq:transport-small}, and combine the Gaussian prefactor using \eqref{eq:element-curvature}. Since $q_0(u+1)=Nu(u+1)/(u+2)$, this proves \eqref{eq:bcarrier}.
\end{proof}
In particular, the proof gives the coarse consequence
\begin{equation}\label{eq:loggrowth}
 \log b_N=N[\log N-2\log\log N+\log2-1+o(1)].
\end{equation}
As explained in Section~\ref{sec:models}, the leading scale is prior and finer logarithmic expressions are already recorded in OEIS. The relative statement \eqref{eq:bcarrier} is stronger than merely exponentiating a logarithmic-scale equivalence. Replacing $u$ by a fixed finite logarithmic expansion generally destroys its relative accuracy because the remaining error is multiplied by the large size in the exponent.
